% capitolo 7: data governance e sviluppo sicuro
\chapter{Data Governance e Sviluppo Sicuro}
\label{cap:sviluppo-sicuro}

\intro{Questo capitolo affronta le tematiche legate alla sicurezza e alla conformità normativa, in linea con i principi della \emph{Privacy by Design}.}\\

\section{Introduzione allo sviluppo sicuro}
L'introduzione di un Data Catalog all'interno dell'infrastruttura aziendale non risponde solo alle esigenze di catalogazione dei dati e di gestione dei metadati, 
ma rappresenta anche uno strumento per garantire la conformità normativa dei dati, l'adesione alle policy aziendali e il rispetto del \gls{gdpr}, nonché un mezzo per supportare il Ciclo di Sviluppo del Software Sicuro.

A tal proposito, Wintech mette a disposizione una serie di documenti e linee guida per la protezione della privacy, la gestione dei dati sensibili e lo sviluppo sicuro del software.

All'interno del contesto aziendale, viene definito il concetto di \emph{Privacy by Design} (PbD), intesa come un approccio olistico concettuale che può essere applicato \emph{end-to-end} all'interno di un'organizzazione, 
includendo le sue tecnologie informatiche, le sue pratiche commerciali, i suoi processi, la progettazione fisica e le infrastrutture di rete. 
Questa definizione si basa su tre concetti fondamentali:
\begin{itemize}
    \item \textbf{Integrazione nativa fin dalla progettazione}: Qualsiasi progetto deve essere realizzato considerando la riservatezza e la protezione dei dati personali sin dalla sua prima fase di progettazione e nell'architettura dei sistemi IT. Non è una misura attuata successivamente a un evento o a un problema sorto.
    \item \textbf{Approccio \emph{user centric}}: Secondo questa impostazione, l'utente deve essere posto al centro del sistema di protezione dei dati personali.
    \item \textbf{Azione proattiva e preventiva}: La PbD è caratterizzata dall'adozione di misure proattive piuttosto che reattive. Il suo scopo primario è anticipare e prevenire gli eventi lesivi per la privacy prima che accadano, invece di attendere che il rischio si materializzi per poi correggerlo.
\end{itemize}

\subsection{Principi della privacy}
Nella definizione delle proprie Linee Guida allo Sviluppo Sicuro, Wintech si ispira allo standard \textbf{ISO/IEC 29100:2011}, il quale raccoglie undici principi fondamentali 
che indirizzano la progettazione, lo sviluppo e l'implementazione dei requisiti di protezione della privacy all'interno dei software. 
Tali principi sono descritti in dettaglio di seguito:

% \begin{longtable}{ p{0.30\textwidth} p{0.64\textwidth} }
\begin{longtable}{ >{\raggedright\arraybackslash}p{0.30\textwidth} p{0.64\textwidth} }
\toprule
\textbf{Principio} & \textbf{Descrizione} \\
\midrule
\endfirsthead

\toprule
\textbf{Principio} & \textbf{Descrizione} \\
\midrule
\endhead

\midrule
\noalign{\vspace{10pt}}
\caption[]{Principi della privacy secondo lo standard ISO/IEC 29100:2011} \\
\endfoot

\bottomrule
\noalign{\vspace{10pt}}
\caption{Principi della privacy secondo lo standard ISO/IEC 29100:2011}
\label{tab:principi-privacy}
\endlastfoot
    \textbf{Consenso e Scelta} & L'interessato deve poter scegliere se acconsentire o meno al trattamento, in modo chiaro e accessibile.\\ \midrule
    \textbf{Legittimità e Specifica delle Finalità} & Il trattamento deve essere conforme alla legge e basato su una base giuridica ammissibile.\\ \midrule
    \textbf{Limitazione della Raccolta} & La raccolta dei dati deve limitarsi a ciò che è strettamente necessario per gli scopi prefissati.\\ \midrule
    \textbf{Minimizzazione dei Dati} & I sistemi ICT devono essere progettati per ridurre al minimo i dati personali elaborati.\\ \midrule
    \textbf{Limitazione dell'Utilizzo, Conservazione e Divulgazione} & I dati vanno usati e conservati solo per scopi specifici, espliciti e legittimi.\\ \midrule
    \textbf{Accuratezza e Qualità} & I dati elaborati devono essere accurati, completi, aggiornati e pertinenti.\\ \midrule
    \textbf{Apertura, Trasparenza e Informativa} & Devono essere fornite informazioni chiare e accessibili sulle procedure di trattamento.\\ \midrule
    \textbf{Partecipazione Individuale e Accesso} & Agli interessati deve essere garantita la possibilità di accedere e rivedere i propri dati.\\ \midrule
    \textbf{Responsabilizzazione} & Le politiche di privacy devono essere documentate, comunicate e assegnate a un responsabile specifico.\\ \midrule
    \textbf{Sicurezza delle Informazioni} & I dati devono essere protetti da rischi come accesso non autorizzato, modifica, divulgazione o perdita.\\ \midrule
    \textbf{Conformità alla Privacy} & Devono essere previste verifiche periodiche per dimostrare il rispetto della tutela della privacy.\\
\end{longtable}

\subsection{Adeguamento al GDPR nel Ciclo di Sviluppo del Software Sicuro}
Il \gls{gdpr} ha un impatto profondo sul \gls{sdlc} per tutte le organizzazioni che progettano e implementano sistemi informativi all'interno dell'Unione Europea. 
I requisiti funzionali e tecnici introdotti dal Regolamento influenzano quasi tutti gli aspetti della progettazione architetturale, a partire 
dai livelli più bassi, come l'architettura dei dati (database) e i livelli di trasporto, per arrivare fino ai livelli applicativi e di presentazione. 
Wintech si impegna a soddisfare tali requisiti fin dalle prime fasi dell'SDLC.

Di seguito vengono mappati e spiegati alcuni dei principali articoli del GDPR e le misure che l'azienda si impegna ad adottare per garantire la \emph{Privacy by Design} e lo Sviluppo Sicuro.
% \begin{longtable}{ p{0.18\textwidth} p{0.37\textwidth} p{0.37\textwidth} }
\begin{longtable}{ >{\raggedright\arraybackslash}p{0.18\textwidth} p{0.37\textwidth} p{0.37\textwidth} }
\toprule
\textbf{Articolo del GDPR} & \textbf{Principio normativo} & \textbf{Impegno Wintech} \\
\midrule
\endfirsthead

\toprule
\textbf{Articolo del GDPR} & \textbf{Principio normativo} & \textbf{Impegno Wintech} \\
\midrule
\endhead

\midrule
\noalign{\vspace{10pt}}
\caption[]{Mappatura dei principali articoli del GDPR rispetto agli impegni aziendali} \\
\endfoot

\bottomrule
\noalign{\vspace{10pt}}
\caption{Mappatura dei principali articoli del GDPR rispetto agli impegni aziendali}
\label{tab:articoli-gdpr}
\endlastfoot
    \textbf{Art. 5 -- Principi applicabili al trattamento dei dati personali} & I dati personali sono raccolti in modo lecito, corretto e trasparente nei confronti dell'interessato e per finalità determinate, esplicite e legittime. Devono essere adeguati, pertinenti e limitati a quanto necessario per le finalità dichiarate; devono essere esatti, aggiornati e conservati per un arco di tempo non superiore al conseguimento delle finalità stesse. & Sotto la direzione del titolare, i dati sono protetti garantendone l'integrità e la riservatezza. Wintech implementa strategie di gestione dei registri elettronici affinché i dati vengano conservati solo per il tempo strettamente necessario; i dati personali scaduti non vengono memorizzati e l'interessato viene informato di tale periodo.\\ \midrule
    \textbf{Art. 13 -- Informazioni da fornire al momento della raccolta} & Il titolare del trattamento deve fornire all'interessato le proprie informazioni di contatto, le finalità e la base giuridica del trattamento, gli eventuali destinatari, il periodo di conservazione, nonché l'esistenza dei diritti di accesso, rettifica, cancellazione e revoca del consenso. & I sistemi sono progettati affinché, al momento della raccolta dei dati, gli attributi del trattamento e le fasi elaborative vengano forniti all'utente per via elettronica o per iscritto, in forma chiara e facilmente comprensibile, garantendo assoluta trasparenza.\\ \midrule
    \textbf{Art. 15 -- Diritto di accesso dell'interessato} & L'interessato ha il diritto di ottenere dal titolare del trattamento la conferma che sia o meno in corso un trattamento di dati personali che lo riguardano e, in tal caso, di ottenerne l'accesso. & Le architetture software devono garantire all'interessato il diritto di accedere ai propri dati e di controllarne il trattamento in qualsiasi momento. Inoltre, l'infrastruttura deve prevedere la possibilità di fornire all'interessato una copia dei propri dati in un formato comunemente utilizzato.\\ \midrule
    \textbf{Art. 16 -- Diritto di rettifica} & L'interessato ha il diritto di ottenere dal titolare del trattamento la rettifica dei dati personali inesatti che lo riguardano senza ingiustificato ritardo e l'integrazione dei dati incompleti. & I sistemi informativi devono consentire all'interessato di ottenere gratuitamente l'aggiornamento dei propri dati in caso di errore. L'infrastruttura deve supportare l'esercizio di tale diritto richiedendo l'aggiornamento direttamente per via elettronica.\\ \midrule
    \textbf{Art. 17 -- Diritto alla cancellazione (Diritto all'oblio)} & L'interessato ha il diritto di ottenere dal titolare la cancellazione dei dati personali che lo riguardano senza ingiustificato ritardo, qualora sussistano i motivi previsti dal Regolamento (es. revoca del consenso o cessazione delle finalità). & I sistemi informativi devono permettere la cancellazione immediata dei dati qualora richiesta. Inoltre, il dipartimento IT si impegna a mantenere inventari aggiornati dei dati interni ed esterni, in modo da sapere esattamente dove tali dati sono conservati da terze parti e poter comunicare tempestivamente l'ordine di cancellazione ad altre organizzazioni IT.\\ \midrule
    \textbf{Art. 25 -- Protezione dei dati fin dalla progettazione} & Il titolare del trattamento mette in atto misure tecniche e organizzative adeguate a garantire che siano trattati solo i dati personali necessari per ogni specifica finalità, integrandovi le necessarie garanzie a tutela dei diritti degli interessati. & L'implementazione della protezione dei dati all'interno del sistema informativo e dell'organizzazione avviene in modo nativo, per progettazione e come impostazione predefinita. Wintech considera questo approccio un requisito necessario nello sviluppo IT.\\ \midrule
    \textbf{Art. 32 -- Sicurezza del trattamento} & Il titolare e il responsabile del trattamento mettono in atto misure tecniche e organizzative per garantire un livello di sicurezza adeguato al rischio, tra cui la pseudonimizzazione e la cifratura dei dati personali. & L'infrastruttura IT di Wintech deve essere progettata avvalendosi di mezzi tecnici e organizzativi rigorosi. Ove possibile, i sistemi software sviluppati devono prevedere sempre l'uso della cifratura dei dati, nonché supportare l'anonimizzazione o la pseudonimizzazione delle informazioni.\\ \midrule
    \textbf{Artt. 33 e 34 -- Notifica e comunicazione di violazione dei dati} & In caso di violazione della sicurezza che comporta la perdita o la divulgazione illecita di dati, il titolare deve notificarla all'autorità di controllo (entro 72 ore) e all'interessato, qualora la violazione presenti un rischio elevato per i suoi diritti. & I sistemi software devono essere dotati di un ambiente tecnico opportuno per individuare, tracciare e valutare quasi immediatamente le violazioni dei dati. Questo è necessario per stabilire in tempo reale se vi sia un rischio elevato e adempiere agli obblighi di comunicazione.\\
\end{longtable}

La gestione di tutti i requisiti normativi sopraelencati richiede una documentazione costante delle fasi di progettazione e un supporto tecnologico adeguato. 
In questo scenario, l'adozione di un Data Catalog aziendale come OpenMetadata diventa una componente architetturale molto importante: automatizzando l'individuazione dei dati sensibili, 
applicando i relativi tag di sicurezza e mappando i flussi informativi, il catalogo fornisce all'azienda uno strumento concreto per garantire la conformità al ciclo di sviluppo sicuro  
e alle direttive imposte dal GDPR, trasformando i vincoli legali in regole tecniche implementate in modo proattivo.


\section{Classificazione dei dati}
Nel contesto della Data Governance, la classificazione dei dati rappresenta il processo logico mediante il quale si associano etichette semantiche (tag) agli asset informativi 
per descriverne in modo standardizzato la natura, l'area di business, l'appartenenza aziendale e il livello di riservatezza. Questa operazione agevola le attività di Data Discovery, 
consentendo agli utenti di cercare e filtrare gli elementi del catalogo, e sta alla base del rispetto delle policy di sicurezza e della conformità normativa, con particolare riferimento al GDPR.

Per configurare correttamente un catalogo dati aziendale, è fondamentale riconoscere quali tipologie di informazioni necessitano di maggiore tutela, come i 
dati sensibili e personali. A tal proposito, l'Articolo 4 del GDPR fornisce la seguente definizione di \textbf{dato personale}~\cite{site:gdpr}:
\begin{quoting}
    «Qualsiasi informazione riguardante una persona fisica identificata o identificabile (\emph{interessato}); si considera identificabile la persona fisica che può essere identificata, direttamente o indirettamente, con particolare riferimento a un identificativo come 
    il nome, un numero di identificazione, dati relativi all'ubicazione, un identificativo online o a uno o più elementi caratteristici della sua identità fisica, fisiologica, genetica, psichica, economica, culturale o sociale.»
\end{quoting}

Le linee guida aziendali di Wintech si adeguano a questa normativa, classificando le informazioni nelle seguenti macro-categorie:
\begin{itemize}
    \item \textbf{Dati identificativi (\gls{pii})}: Dati che consentono di identificare direttamente o indirettamente una persona.
    \item \textbf{Dati particolari/sensibili (Art. 9 GDPR)}: Dati personali che necessitano di una maggiore protezione poiché il loro trattamento potrebbe creare rischi significativi ai diritti o alle libertà fondamentali dell'interessato; si tratta per esempio di: origine razziale o etnica, opinione politica, ideologie religiose o filosofiche, appartenenza sindacale, dati genetici o biometrici, dati relativi alla salute o alla vita sessuale.
    \item \textbf{Dati giudiziari}: Dati relativi alle condanne penali e ai reati o a connesse misure di sicurezza; la gestione di queste informazioni è consentita sotto il controllo di un'autorità pubblica.
\end{itemize}

\noindent Viene fatta anche una distinzione tra i seguenti concetti:
\begin{itemize}
    \item \textbf{Dati pseudonimizzati}: Dati che vengono sostituiti da pseudonimi, in cui l'assegnazione degli pseudonimi è tale da non poter essere invertita e gli attributi collegati allo pseudonimo non sono sufficienti per identificare l'interessato. Tuttavia, i dati pseudonimizzati possono essere violati e l'interessato può essere identificato, pertanto sono soggetti al GDPR.
    \item \textbf{Dati anonimizzati}: Informazioni modificate in modo irreversibile, al punto che il titolare del trattamento non può identificare l'interessato. I dati resi anonimi non sono più considerati personali e quindi escono dal perimetro del GDPR.
    \item \textbf{Metadati}: I dati personali possono essere memorizzati in un sistema ICT in modo tale da non essere facilmente visibili all'utente del sistema. L'interessato deve essere a conoscenza dell'esistenza di tali dati sotto forma di metadati o del loro trattamento per tale scopo.
    \item \textbf{Dati non richiesti}: Informazioni personali fornite volontariamente dall'utente anche quando non richieste dal trattamento. Il rischio di raccogliere informazioni personali indesiderate può essere ridotto considerando le misure di tutela della privacy al momento della progettazione del software.
\end{itemize}

\noindent Data la specificità di queste categorie, una gestione manuale efficace della classificazione e del controllo dei dati personali risulterebbe molto complessa. 
In questo contesto, la funzionalità di auto-classificazione di un Data Catalog come OpenMetadata assume un ruolo chiave: la piattaforma scansiona gli asset di dati 
alla ricerca di queste categorie, sfruttando motori di Intelligenza Artificiale o espressioni regolari.

\subsection{Tassonomia dei tag}
OpenMetadata organizza i tag all'interno di categorie logiche centralizzate, denominate \emph{Classifications}. Per garantire una corretta gestione dei 
metadati e tradurre i requisiti normativi in regole tecniche, la piattaforma fornisce alcune classificazioni di sistema predefinite. Tali categorie non possono essere 
eliminate, ma guidano l'assegnazione corretta delle etichette:
\begin{itemize}
    \item \textbf{PII}: Classificazione standard concepita per l'identificazione e la mappatura delle informazioni personali, fondamentale per la conformità al GDPR. Include il tag \texttt{PII.NonSensitive}, utilizzato per dati identificativi generici e a basso rischio (come nome, cognome o e-mail aziendale), e il tag \texttt{PII.Sensitive}, obbligatorio per gli identificatori critici (come codici fiscali, password o coordinate bancarie). L'assegnazione di quest'ultimo innesca automaticamente logiche di anonimizzazione per gli utenti sprovvisti dei privilegi giusti.
    \item \textbf{PersonalData}: Classificazione complementare per mappare i dati riferibili a persone fisiche, suddivisa in \texttt{PersonalData.Personal} e \texttt{PersonalData.SpecialCategory}. Quest'ultimo tag è destinato ai dati particolari (Art. 9 del GDPR, come dati biometrici o sanitari) e ai dati giudiziari, imponendo un approccio di massima chiusura che ne limita l'accesso esclusivamente alle autorità o ai ruoli aziendali espressamente autorizzati.
    \item \textbf{\emph{Tier} (Livello)}: Utilizzata per suddividere gli asset in base alla loro rilevanza strategica e architetturale per l'organizzazione. Comprende cinque livelli gerarchici standard, da \texttt{Tier.Tier1} (assegnato ad asset contenenti dati critici e insostituibili per il business) fino a \texttt{Tier.Tier5} (associato a informazioni marginali o secondarie).
    \item \textbf{\emph{Certification} (Certificazione)}: Impiegata all'interno delle procedure di Data Quality per qualificare visivamente l'affidabilità di un dataset. Questa tassonomia si articola nei tag gerarchici \texttt{Bronze}, \texttt{Silver} e \texttt{Gold}, permettendo agli analisti di Business Intelligence di verificare all'istante se un dato è stato validato e approvato prima di importarlo nei report.
\end{itemize}

\begin{figure}[H]
    \centering
    \begin{subfigure}[b]{0.9\textwidth}
        \centering
        \includegraphics[width=\textwidth]{schermate/tag_pii.png}
        \caption{Tassonomia PII}
    \end{subfigure}
    
    \vspace{0.4cm}
    
    \begin{subfigure}[b]{0.9\textwidth}
        \centering
        \includegraphics[width=\textwidth]{schermate/tag_tier.png}
        \caption{Tassonomia Tier}
    \end{subfigure}
    
    \caption{Pannelli di gestione delle classificazioni di sistema}
    \label{fig:tassonomie_sistema}
\end{figure}

\noindent Oltre alle tassonomie predefinite, la piattaforma supporta la creazione di classificazioni personalizzate, adattabili alle specifiche policy interne 
di Wintech, per esempio per contrassegnare dati con un tag che ne indica il livello di «confidenzialità».

Ogni classificazione può essere configurata attivando l'opzione di mutua esclusione. Se abilitata, questa proprietà garantisce l'integrità logica della classificazione, 
impedendo che a un singolo asset vengano associati contemporaneamente due tag contrastanti appartenenti alla medesima tassonomia.

\begin{figure}[H]
    \centering
    \includegraphics[width=0.75\textwidth]{schermate/tag_personalizzato.png}
    \caption[Esempio di classificazione personalizzata]{Esempio di classificazione personalizzata \texttt{DataClassification}, configurata con l'opzione di mutua esclusione (\emph{Mutually Exclusive}) attivata}
    \label{fig:tag_personalizzato}
\end{figure}

\subsection{Classificazione automatica}
L'\emph{Auto-Classification} è un workflow di OpenMetadata che rileva e classifica automaticamente i dati sensibili nelle colonne dei database, eliminando la necessità 
di tagging manuale. Questa procedura, eseguita in background e integrabile all'interno delle \emph{pipeline} di sincronizzazione programmata, sfrutta l'operato 
dell'\emph{AutoClassification Agent} per scansionare le sorgenti collegate, applicando o suggerendo etichette come \texttt{PII.Sensitive} e \texttt{PII.NonSensitive}.

Il motore di \emph{auto-classificazione} utilizza due approcci complementari: 
\begin{enumerate}
    \item \textbf{\emph{Column Name Scanner}}: Esegue una scansione strutturale rapida dei metadati, confrontando i nomi fisici delle colonne del database con un elenco di espressioni regolari (regex) predefinite o parole chiave note (identificando, ad esempio, colonne nominate \texttt{email}, \texttt{phone} o \texttt{fiscal\_code}).
    \item \textbf{\emph{Named Entity Recognition} (NER)}: Qualora sia stata abilitata la raccolta di dati campione (\emph{Sample Data}) durante la fase di \emph{ingestion}, il motore NLP analizza il contenuto testuale reale dei record estratti per dedurne semanticamente la natura, identificando le informazioni personali anche in presenza di colonne con nomenclatura generica o non strutturata.
\end{enumerate}

\begin{figure}[H]
    \centering
    \includegraphics[width=1\textwidth]{schermate/dati_sensibili.png}
    \caption[Tabella con tag PII applicati automaticamente]{Esempio di una tabella in cui il motore di auto-classificazione ha applicato automaticamente il tag \texttt{PII.Sensitive} alle colonne contenenti dati personali}
    \label{fig:dati_sensibili}
\end{figure}

\subsubsection{Funzionamento dell'algoritmo di NER}
Il modulo di NER si basa sul \emph{framework} \emph{open-source} \textbf{Microsoft Presidio}, uno strumento specializzato nella rilevazione e anonimizzazione 
dei dati personali che si appoggia a librerie NLP (come spaCy o Stanza). 

In OpenMetadata, Presidio opera come pacchetto Python eseguito direttamente all'interno dell'ambiente di Data Ingestion, gestendo il flusso dei dati in modo \emph{stateless}. 
I dati campione prelevati dal database vengono caricati temporaneamente nella RAM del Container ed elaborati dal modello, il quale restituisce un file in formato JSON 
contenente le informazioni rilevate. Al termine della scansione, il campione viene immediatamente distrutto dalla memoria e l'agente trasmette al server centrale 
esclusivamente il metadato risultante (il tag di classificazione da applicare).

\begin{figure}[H]
    \centering
    \includegraphics[width=0.75\textwidth]{autoclassificazione.png}
    \vspace{0.2cm}
    \caption[Rappresentazione del flusso di auto-classificazione]{Architettura dell'auto-classificazione: i dati reali vengono elaborati e distrutti localmente, trasmettendo al server centrale solo i metadati. \textit{Nota: La presente infografica è stata generata con l'ausilio di un sistema di Intelligenza Artificiale}}
    \label{fig:flusso_autoclassificazione}
\end{figure}

Per rispettare il GDPR e i vincoli di sicurezza imposti dall'infrastruttura di Wintech, l'intero processo avviene in locale e offline. A differenza degli assistenti basati 
su Cloud, i modelli linguistici necessari per l'analisi vengono scaricati una sola volta durante la fase iniziale di compilazione e installazione del Container. 
A riprova di questo, su reti aziendali dotate di rigide policy di sicurezza, il Container può restituire l'errore \texttt{SSL: CERTIFICATE\_VERIFY\_FAILED} 
al primo avvio, venendo bloccato dal firewall aziendale durante il tentativo di download del modello; una volta scaricato manualmente, il sistema opera totalmente 
offline garantendo che nessuna informazione sensibile venga trasmessa all'esterno della rete aziendale.

Infine, l'architettura basata su Presidio garantisce una notevole scalabilità tramite i cosiddetti \emph{Custom Recognizers}. Qualora l'azienda necessiti di intercettare 
formati di dati specifici o interni (per esempio la sintassi proprietaria di un ID Cliente), è sufficiente sviluppare e integrare un riconoscitore personalizzato 
all'interno della \emph{pipeline} Python, estendendo le capacità di rilevamento senza alterare i protocolli di sicurezza.


\subsection{Limiti dell'automazione e classificazione manuale}
Durante i test di implementazione, è emerso che la funzionalità di classificazione automatica non è esente da imprecisioni contestuali e presenta alcune limitazioni architetturali. 

In primo luogo, l'efficacia del rilevamento automatico sui database aziendali dipende dall'ottimizzazione di parametri tecnici specifici:
\begin{itemize}
    \item \textbf{Propagazione dei prerequisiti logici}: l'esecuzione dell'\emph{AutoClassification Agent} è strettamente subordinata all'avvenuto completamento delle attività del \emph{Profiler Agent}. Poiché il motore NLP necessita del contenuto testuale reale dei record (\emph{Sample Data}) per applicare la Named Entity Recognition, è indispensabile pianificare l'esecuzione delle \emph{pipeline} garantendo che il profilo dei dati venga generato prima della scansione di classificazione;
    \item \textbf{Regolazione del livello di fiducia (\emph{confidence level})}: per bilanciare la precisione del rilevamento, il sistema consente di calibrare una soglia minima di confidenza. Questo parametro è fondamentale per evitare sia l'insorgere di falsi positivi (dati innocui scambiati per sensibili) sia di falsi negativi (dati sensibili non rilevati).
\end{itemize}

In secondo luogo, è stato rilevato un importante vincolo relativo alla versione \emph{open-source} di OpenMetadata: il modulo è progettato per operare una classificazione 
prevalentemente binaria a livello di GDPR, etichettando i campi sensibili in modo generico come \texttt{PII.Sensitive} o \texttt{PII.NonSensitive}. La capacità di associare 
automaticamente tag più granulari (come \texttt{Address} o \texttt{Name}) è una funzionalità esclusiva della distribuzione commerciale (\emph{Collate}) e non è direttamente 
fruibile nell'installazione locale standard. 

A fronte di questi limiti tecnici e per garantire la massima conformità alle policy aziendali, è indispensabile una fase di classificazione e revisione manuale. 
In questo passaggio, figure qualificate come i \emph{data steward} o i proprietari del dato (\emph{data owners}) intervengono per validare, correggere o raffinare le 
etichette suggerite dalla piattaforma. Questa attività permette anche di arricchire i metadati dell'asset: l'utente può inserire descrizioni contestuali dettagliate, 
definire l'appartenenza a specifici domini aziendali e associare i campi ai termini formalizzati nel glossario.

\subsection{Integrazione dell'AI Generativa tramite MCP}
A completamento delle analisi architetturali svolte durante lo stage, si è esplorata la possibilità di integrare strumenti di Intelligenza Artificiale generativa 
in Cloud (come ChatGPT o Claude) per abilitare la ricerca semantica e l'interrogazione del catalogo in linguaggio naturale. Questo è possibile grazie all'integrazione 
nativa all'interno di OpenMetadata, disponibile anche nella versione \emph{open-source}, di un server basato sullo standard \textbf{\gls{mcp}}. 

Il protocollo MCP fa da interfaccia di comunicazione sicura tra il catalogo dati e gli assistenti AI esterni. Condividendo esclusivamente i metadati, 
questa soluzione risponde ai requisiti di sicurezza attraverso i seguenti meccanismi:

\begin{itemize}
    \item \textbf{Totale isolamento dei dati reali}: L'assistente AI non stabilisce alcuna connessione diretta con i database fisici dell'azienda, precludendo 
    l'accesso ai dati in chiaro. Il modello interagisce unicamente con il server MCP di OpenMetadata, il quale trasmette all'esterno esclusivamente metadati 
    strutturali (schema delle tabelle, nomi dei campi, definizioni del glossario, tag di classificazione applicati). Le informazioni effettive contenute nei 
    database non vengono mai elaborate dall'AI e non lasciano il perimetro aziendale, garantendo la conformità alle direttive sulla sicurezza.
    \item \textbf{Preservazione del controllo degli accessi}: L'utilizzo dell'AI tramite server MCP è subordinato al rispetto delle policy di Data Governance, impedendo 
    all'Intelligenza Artificiale di bypassare i sistemi di sicurezza interni. La connessione al server richiede un'autenticazione sicura: l'assistente AI eredita 
    direttamente i permessi dell'utente che ha effettuato l'accesso, seguendo il modello di controllo degli accessi basato sui ruoli (RBAC) configurato sulla piattaforma. 
    Se l'account dell'utente che interroga l'AI non dispone delle autorizzazioni necessarie per visualizzare una determinata area del catalogo, il server MCP bloccherà 
    la richiesta alla fonte, impedendo all'AI di accedere ai relativi metadati e di includerli nella risposta.
\end{itemize}


\section{Mascheramento dei dati}
A differenza di quanto ipotizzato in fase di analisi dei requisiti, la piattaforma non dispone di impostazioni personalizzabili di anonimizzazione o pseudonimizzazione dei dati. 
Tuttavia, OpenMetadata integra un meccanismo automatico di mascheramento totale (anonimizzazione) strettamente dipendente dal modello di accesso basato sui ruoli.

Quando una colonna o una tabella viene classificata con il tag \texttt{PII.Sensitive}, i relativi \emph{Sample Data} vengono oscurati e sostituiti visivamente da asterischi 
(\texttt{******}) all'interno dell'interfaccia. Questo blocco visivo viene applicato in automatico a tutti gli utenti standard, garantendo invece la visualizzazione dei valori 
in chiaro esclusivamente agli amministratori di sistema (\emph{Admin}) e ai proprietari (\emph{owner}) del singolo asset.

\begin{figure}[H]
    \centering
    \includegraphics[width=1\textwidth]{schermate/anonimizzazione.png}
    \vspace{0.2cm}
    \caption[Applicazione dell'anonimizzazione nell'interfaccia]{Applicazione pratica dell'anonimizzazione: i \emph{Sample Data} classificati come sensibili vengono oscurati per gli utenti non autorizzati}
    \label{fig:anonimizzazione}
\end{figure}

\subsection{Limiti architetturali}
Durante l'analisi e il collaudo della funzionalità, sono state riscontrate alcune limitazioni:
\begin{itemize}
    \item \textbf{Assenza di pseudonimizzazione}: Non è possibile scegliere tra anonimizzazione e pseudonimizzazione. Il comportamento del meccanismo di mascheramento è predefinito e si applica in modo uniforme (mascheramento totale) a tutti gli asset etichettati come \texttt{PII.Sensitive}.
    \item \textbf{Rigidità configurativa}: Non è permesso configurare autonomamente regole di mascheramento basate sull'assegnazione di tag di classificazione personalizzati.
    \item \textbf{Mancata propagazione esterna}: Non sono previste funzionalità per propagare le direttive di mascheramento ai sistemi esterni a valle (come, ad esempio, gli strumenti di Business Intelligence).
    \item \textbf{Dipendenza dal modello RBAC}: Il meccanismo si basa interamente sul sistema di \emph{ownership} e ruoli. Una rigorosa configurazione dei permessi rappresenta quindi un prerequisito architetturale indispensabile affinché la protezione dei dati sensibili sia garantita.
\end{itemize}

Nonostante questi vincoli tecnici, OpenMetadata offre una funzionalità di base per l'anonimizzazione dei dati che risulta comunque sufficiente per soddisfare i requisiti normativi e le esigenze aziendali.


\section{Gestione degli accessi}
L'efficacia della Data Governance e la conformità al GDPR richiedono una gestione degli accessi sicura, granulare e centralizzata. L'adozione di un Data Catalog 
aziendale deve garantire il rispetto del \emph{Principle of Least Privilege}, fungendo da unico punto di controllo per l'accesso ai metadati e proteggendo in particolare le informazioni sensibili.

Per rispondere a questa esigenza, OpenMetadata utilizza un modello di controllo degli accessi basato sui ruoli (RBAC), combinato con una fedele mappatura 
gerarchica dell'organigramma aziendale, possibile grazie alla definizione di utenti e team.

\subsection{Mappatura dell'organigramma: utenti e team}
Il primo livello di sicurezza riguarda la gestione delle identità. Dalla scheda di amministrazione è possibile configurare i singoli utenti (\emph{User}) e 
raggrupparli in gruppi di lavoro (\emph{Team}). Questa struttura permette di replicare fedelmente l'organigramma di Wintech all'interno del sistema. 
OpenMetadata permette di organizzare l'albero aziendale secondo cinque livelli logici (\emph{teamType}):
\begin{itemize}
    \item \textbf{\emph{Organization}}: È il nodo radice della gerarchia. Contiene tutti i dipendenti e i sottogruppi dell'azienda e non possiede nodi genitori.
    \item \textbf{\emph{BusinessUnit}}: Rappresenta le divisioni di primo livello dell'azienda (es. Direzione, Ricerca e Sviluppo). Può dipendere esclusivamente dall'\emph{Organization} o da un'altra \emph{BusinessUnit}.
    \item \textbf{\emph{Division}}: È una partizione operativa intermedia. Dipende gerarchicamente da un solo nodo superiore (\emph{Organization}, \emph{BusinessUnit} o un'altra \emph{Division}).
    \item \textbf{\emph{Department}}: Individua aree di competenza specialistiche. A differenza degli altri livelli, può ereditare proprietà da più nodi genitori, offrendo maggiore flessibilità organizzativa.
    \item \textbf{\emph{Group}}: Rappresenta il nodo foglia della gerarchia. Non può contenere sottogruppi, ma è formato solo da singoli \emph{User}. Nel Data Catalog, solo i team di tipo \emph{Group} possono detenere formalmente la proprietà (\emph{Ownership}) degli asset.
\end{itemize}

L'inserimento degli utenti può avvenire manualmente, ma l'infrastruttura supporta anche la centralizzazione degli accessi tramite l'integrazione con protocolli di 
Single Sign-On (SSO), forniti da provider come Microsoft Entra ID, Google Workspace e Okta. 
Questa configurazione abilita il provisioning automatico (\emph{Auto-Provisioning}). Al momento del primo accesso di un utente tramite credenziali aziendali, 
la piattaforma interroga il provider e crea in automatico il relativo profilo nel catalogo, riducendo l'onere amministrativo per la gestione del personale.

\begin{figure}[H]
    \centering
    \begin{subfigure}[b]{1\textwidth}
        \centering
        \includegraphics[width=\textwidth]{schermate/user.png}
        \caption{Pannello di amministrazione dei singoli utenti (\emph{User})}
    \end{subfigure}
    
    \vspace{0.4cm}
    
    \begin{subfigure}[b]{1\textwidth}
        \centering
        \includegraphics[width=\textwidth]{schermate/team.png}
        \caption{Mappatura gerarchica dell'organigramma tramite i \emph{Team}}
    \end{subfigure}
    
    \caption[Pannelli di gestione degli utenti e dei team]{Interfacce per la gestione delle identità in OpenMetadata: visualizzazione dell'elenco utenti (sopra) e dell'organizzazione gerarchica in dipartimenti o \emph{BusinessUnit} (sotto)}
    \label{fig:gestione_utenti_team}
\end{figure}

\subsection{L'architettura Role-Based Access Control}
Una volta definita la struttura degli utenti e dei team, la piattaforma permette di definire autorizzazioni e permessi attraverso il sistema \gls{rbac}.
Dal punto di vista architetturale, il motore di controllo degli accessi si basa su tre componenti interconnessi:
\begin{enumerate}
    \item \textbf{Rules (Regole)}: Sono l'elemento base della sicurezza. Definiscono se una specifica azione (es. modificare tag, visualizzare dati campione) è permessa o negata su un certo tipo di risorsa.
    \item \textbf{Policies (Politiche)}: Sono raggruppamenti logici di regole. Vengono create per rispondere a specifiche esigenze di sicurezza aziendale (es. policy per la qualità dei dati, policy per la gestione dei domini).
    \item \textbf{Roles (Ruoli)}: Associano una o più politiche ai profili utente. Assegnando un ruolo a un utente o a un intero team, le relative autorizzazioni vengono propagate istantaneamente a tutti i membri.
\end{enumerate}

\subsection{Logica di valutazione dei permessi}
Quando un utente tenta di eseguire un'operazione su un asset, il motore di OpenMetadata valuta la richiesta. La decisione di autorizzare o bloccare l'azione segue 
quattro principi fondamentali:
\begin{itemize}
    \item \textbf{Permessi di base (Default)}: Ogni utente registrato ottiene i privilegi minimi associati al ruolo base (\emph{Data Consumer Role}) e alla policy globale (\emph{OrganizationPolicy}). Questo garantisce che chiunque possa cercare e leggere i dati (Data Discovery), escludendo però le operazioni di modifica.
    \item \textbf{Ereditarietà dai team}: I permessi sono cumulativi. Il sistema calcola l'unione logica tra i privilegi assegnati direttamente al profilo e quelli ereditati da tutti i team a cui l'utente appartiene.
    \item \textbf{Eccezione di Proprietà (\emph{Ownership})}: È una deroga fondamentale alle regole standard. Se un utente o un team è designato come proprietario di un asset, ne ottiene automaticamente la gestione completa, superando i limiti del proprio ruolo ordinario.
    \item \textbf{Priorità dei divieti (\emph{Deny over Allow})}: È la regola di sicurezza che previene i conflitti. Le regole all'interno delle policy possono consentire (\emph{Allow}) o vietare (\emph{Deny}) un'azione. Se un utente eredita un permesso abilitante, ma appartiene a un team con una regola restrittiva esplicita (\emph{Deny}) per la stessa azione, il divieto prevale sempre e l'operazione viene bloccata.
\end{itemize}

\begin{figure}[H]
    \centering
    \begin{subfigure}[b]{1\textwidth}
        \centering
        \includegraphics[width=\textwidth]{schermate/ruolo_datagovmanager.png}
        \caption{Assegnazione di una \emph{policy} di sicurezza allo specifico ruolo di \emph{Data Governance Manager}}
    \end{subfigure}
    
    \vspace{0.4cm}
    
    \begin{subfigure}[b]{1\textwidth}
        \centering
        \includegraphics[width=\textwidth]{schermate/policy_gdpr.png}
        \caption{Dettaglio della \emph{policy}: regola con effetto di divieto (\emph{Deny}) basata sui tag}
    \end{subfigure}
    
    \caption[Esempio di architettura RBAC per la conformità GDPR]{Applicazione pratica dell'architettura RBAC: il Ruolo (sopra) eredita la \emph{policy} dedicata al rispetto del GDPR, la quale (sotto) definisce una regola di divieto (\emph{Deny}) per impedire la visualizzazione in chiaro dei dati etichettati come sensibili}
    \label{fig:rbac_gdpr}
\end{figure}


\section{Glossario aziendale}
Un \textbf{Glossario aziendale} (\emph{Business Glossary}) è un vocabolario centralizzato che raccoglie, definisce e standardizza i concetti e la terminologia 
rilevante all'interno dell'organizzazione. Rappresenta uno strumento di Data Governance fondamentale, a cui ogni utente può fare riferimento per risolvere ambiguità 
semantiche e comprendere appieno il significato del patrimonio informativo.

L'obiettivo principale di questo strumento è stabilire un linguaggio aziendale comune, colmando la distanza comunicativa tra i diversi reparti 
attraverso la definizione di un significato comune per ogni termine. All'interno di OpenMetadata, i concetti formalizzati nel glossario assumono il ruolo di veri e propri 
metadati: possono essere infatti assegnati agli asset per descrivere e categorizzare ulteriormente tabelle e dataset.

Un'organizzazione può definire molteplici glossari indipendenti, strutturandoli in base agli specifici domini aziendali (ad esempio Amministrazione, Risorse Umane o Vendite). 
I termini all'interno di ciascun glossario possono essere organizzati secondo una gerarchia logica e possono essere arricchiti con una serie di proprietà avanzate:
\begin{itemize}
    \item \textbf{Contesto semantico}: descrizioni testuali, sinonimi, riferimenti a link esterni e associazioni con altri termini correlati;
    \item \textbf{Propagazione delle classificazioni}: possibilità di associare tag di classificazione direttamente al termine. Quando il concetto viene applicato a un dataset, quest'ultimo eredita automaticamente anche i tag associati al termine stesso;
    \item \textbf{Vincoli di assegnazione}: funzionalità di mutua esclusione (\emph{Mutually Exclusive}), che impedisce di assegnare simultaneamente due termini concettualmente in conflitto al medesimo asset;
    \item \textbf{Validazione e ownership}: designazione di specifici utenti revisori (\emph{reviewer}), incaricati di validare i termini, approvare le modifiche e garantire l'accuratezza del glossario nel tempo.
\end{itemize}

\begin{figure}[H]
    \centering
    \includegraphics[width=1\textwidth]{schermate/glossario_wintech.png}
    \vspace{0.2cm}
    \caption{Interfaccia principale del Glossario aziendale}
    \label{fig:glossario_generale}
\end{figure}

\begin{figure}[H]
    \centering
    \includegraphics[width=1\textwidth]{schermate/glossario_termini.png}
    \vspace{0.2cm}
    \caption{Struttura gerarchica dei termini nel Glossario}
    \label{fig:glossario_gerarchia}
\end{figure}

\section{Domini e data products}
OpenMetadata permette di organizzare gli asset di dati in gruppi logici, definiti Domini (\emph{domain}), per strutturare il patrimonio informativo in modo più efficace. 
Un Dominio consente di categorizzare e gestire asset, glossari, team e altre entità sotto una struttura logica unificata, supportando un'architettura decentralizzata in 
cui ogni area è gestita in modo indipendente con un'\emph{ownership} distribuita.

\subsection{Valore aziendale e Data Product}
L'introduzione dei Domini permette di decentralizzare la gestione dei dati: delegando la responsabilità ai singoli dipartimenti (come Risorse Umane o Marketing), 
la cura e la documentazione del dato vengono affidate a chi ne possiede le reali competenze. Questo approccio elimina i colli di bottiglia e agevola l'applicazione 
delle policy di sicurezza.

All'interno di ciascun Dominio, gli asset possono essere ulteriormente raggruppati in \textbf{Data Product}. Un Data Product è semplicemente un sottoinsieme di 
dati organizzato per rispondere a un'esigenza specifica (ad esempio, raggruppare solo i dati relativi alle campagne pubblicitarie). Questa suddivisione trasforma 
i dati grezzi in un vero e proprio «prodotto» pronto all'uso, facilmente consultabile dagli utenti o analizzabile tramite strumenti di Business Intelligence.

\begin{figure}[H]
    \centering
    \begin{subfigure}[b]{1\textwidth}
        \centering
        \includegraphics[width=\textwidth]{schermate/domini.png}
        \caption{Vista generale dei Domini configurati nella piattaforma}
    \end{subfigure}
    
    \vspace{0.4cm}
    
    \begin{subfigure}[b]{1\textwidth}
        \centering
        \includegraphics[width=\textwidth]{schermate/data_products.png}
        \caption{Dettaglio dei Data Product (\emph{Advertising} e \emph{Social Media Management}) nel dominio Marketing}
    \end{subfigure}
    
    \caption{Organizzazione logica in Domini e Data Product}
    \label{fig:domini_data_products}
\end{figure}